\documentclass{article}

\setlength{\parindent}{0cm}

\begin{document}
    Um Schreibarbeit und Diagonalmatrizen zu vermeiden sei:

    $$y_i^p = p_i y_i, \quad x_{ij}^p = p_i x_{ij} \quad \text{etc.} $$

\section{Demand Side}
    E gilt:
    $$y_i = \sum_j x_{ij} + c_i$$
    Wobei $y_i$ die produzierte Menge in Sektor $i$ ist, $x_{ij}$ die Menge an Intermediärgütern von $i$ nach $j$ und $c_i$ der final demand ist.

    Diese Gleichung gilt auch nominell:

    $$y_i^p = \sum_j x_{ij}^p + c_i^p \left( = \sum_j p_i x_{ij} + p_i c_i \right)$$

    \subsection{Mengen}
    Definieren wir die technischen Koeffizienten für die Menge wollen wir dass gilt:
    $a_{ij} y_j = x_{ij}$, also
    
    $$a_{ij} = \frac{x_{ij}}{y_j}$$

    Durch einfaches einsetzen folgt also:

    y_i = \sum_{j} a_{ij} y_j + c_i$

    oder eben 
    $$y = A y + c$$

    \subsection{Preise}

    Wie schon bei den Mengen wollen wir $\omega$ so definieren, dass gilt

    $\omega_{ij} y_j^p = x_{ij}^p$

    Also $$\omega_{ij} = \frac{x_{ij}^p}{y_j^p} = \frac{p_i x_{ij}}{p_j y_j}.$$

    Es gilt also wieder:
    $$y_i^p = \sum_j \omega_{ij} y_j^p + c_i^p.$$
    
    \section{Supply side}

    Es gilt;
    $$y_i^p = \sum x_{ji}^p + v_i^p,$$
    wobei $v_i$ die Value-Added in Sektor $i$ ist

    Wir wollen wieder einen technischen Koeffizienten $\psi$ einführen, sodass

    $$\psi_{ij} y_j^p = x_{ji}^p$$
    gilt.
    Um zu verdeutlichen, dass das jetzteinmal nicht das gleiche wie $\omega$ ist habe ich $\psi$ gewählt.
    Also $$\psi_{ij} = \frac{x_{ji}^p}{y_j^p} = \frac{x_{ji}}{y_j}$$

    Und damit 

    $$y_i^p = \sum_j \psi_{ij} y_j^p + v_i^p$$
    
    
 \end{document}